\documentclass[10pt]{article}
\usepackage[margin=20mm]{geometry}
\usepackage{amsmath,amssymb,mathtools,fontspec,unicode-math,xurl,hyperref}
\IfFontExistsTF{DejaVu Serif}{\setmainfont{DejaVu Serif}}{\setmainfont{Latin Modern Roman}}
\setmathfont{latinmodern-math.otf}
\setlength{\parindent}{0pt}\setlength{\parskip}{6pt}\sloppy
\title{Technical Verification of Murmurations}
\author{AI-assisted technical review}
\date{}
\begin{document}
\maketitle
Abstract: A continuation review of the acquired journal version verifies a wrong leading normalization in Theorem 3 and independently proves that the remainder asserted in Lemma 7.3 is incompatible with RH. The latter is a known issue: a later preprint identifies the obstruction and supplies a replacement proof of the sharp-average conclusion. A cutoff error in the derivation of Theorem 1 is also verified and repaired; its effect within the stated parameter range warrants a Major assessment of that printed proof step. The report separates defects in the reviewed version from the survival of the principal murmuration phenomenon and from externally supplied repairs.\par
\section*{Identity, recovery basis, and inspected scope}
The reviewed paper is Nina Zubrilina, Murmurations, Inventiones mathematicae 241 (2025), 627–680, DOI 10.1007/s00222-025-01347-8 (https://doi.org/10.1007/s00222-025-01347-8). The primary input is the acquired 54-page publisher Version of Record. PDF p. means its one-based PDF page; printed page equals PDF page plus 626. The separate companion arXiv:2310.07681v2 is used only for a static comparison of the cited formulas. Whole-document equivalence is not asserted.\par
This continuation recovers the substantive report and mathematical checks produced by the predecessor recovery. Its analytic arguments, exact checks, source extraction, and numerical output are retained with their original attribution. The present recovery checks the cited source correspondence and completes the missing verification record; it does not repeat a whole-paper review.\par
The substantive scope comprises the definitions and main statements in Theorems 1–3 and 5–6 and Corollary 4; the displayed error estimate leading to Proposition 3.1; the Poisson/Bessel reduction in Section 5; and Lemma 7.3 with its use in Section 7.3. Section 8’s computational sign argument was examined for its evidence requirements. The detailed class-number analysis throughout Section 3.3, every intermediate lemma, and the complete proofs in the later repair preprint have not been independently re-proved. No claim of whole-paper correctness is made.\par
\section*{Definitions and operative assumptions}
Equation reference: eq:qbound. Fix an even integer \(k\ge2\). The density \(\mathcal M_k(y)\), for \(y>0\), is defined in Theorem 1. Its weights in Theorem 3 are \[A_p=p^4-2p^2-p+1,\qquad
 Q(d)=\mu^2(d)\prod_{p\mid d}\frac{p^2}{A_p}.\] The constants \(\alpha,\beta,\gamma\) have the paper’s normalization. In particular \(\alpha>0\), \(Q(1)=1\), \(Q(d)\ge0\), and \[
 Q(d)\ll d^{-2},\] because \(\prod_p p^4/A_p\) converges: every factor is positive and \(1+O(p^{-2})\). None of the arguments below changes the normalization of the standard Bessel function \(J_\nu\).\par
For the conditional argument, RH has exactly the meaning used in Lemma 7.3: the nontrivial zeros of the Riemann zeta function have real part \(1/2\). Theorem 6 additionally assumes fixed \(c>1\) and even \(k\ge6\). For Proposition 3.1, \(P>2\) is prime, \(X\to\infty\), and \(Y=o(X)\); the decisive test below takes \(X=P\) along primes and \(Y=X^{3/4}\), within the actual parameter range of Theorem 1.\par
\section*{Supported findings}
Problem 1 (Incorrect leading normalization in Theorem 3).  The stated asymptotic coefficient differs from the one implied by the exact Bessel representation by the factor \(2\sqrt{\pi}\). Its effect is a leading-order discrepancy, not an error absorbed by the stated \(O(1)\).\par
\subsection*{Source and decisive calculation}
Equation reference: eq:exact. Theorem 3, PDF p. 5 (printed p. 631), states \[
 \mathcal M_k(y)=\alpha\sqrt y\sum_{d,s\ge1}\frac{Q(d)}s
 J_{k-1}\left(\frac{4\pi s\sqrt y}{d}\right).\] Its next display has the coefficient \(\alpha(-1)^{k/2-1}\sqrt{2/\pi}\) in front of the universal cosine series times \(y^{1/4}\). Section 5, PDF pp. 42–43, derives [eq:exact] (\#eq:exact); its last two displays also introduce a sign discrepancy.\par
Equation reference: eq:correct. Set \(T=\sqrt y\), \(\epsilon_k=(-1)^{k/2-1}\), and \[F(T)=\sum_{d,s\ge1}\frac{Q(d)\sqrt d}{s^{3/2}}
             \cos\left(\frac{4\pi sT}{d}-\frac{3\pi}{4}\right).\] The standard fixed-order expansion [dlmf] gives \[J_{k-1}(z)=\epsilon_k\sqrt{\frac2{\pi z}}\,
              \cos(z-3\pi/4)+O_k(z^{-3/2}),\qquad z\to+\infty.\] The phase follows by subtracting \((k-1)\pi/2+\pi/4\). Substitution into one term of [eq:exact] (\#eq:exact) gives \[\alpha T\,\frac{Q(d)}s\sqrt{\frac2{\pi(4\pi sT/d)}}
 =\frac{\alpha}{\pi\sqrt2}\,T^{1/2}
                 \frac{Q(d)\sqrt d}{s^{3/2}}.\] Consequently the corrected formula is \[
 \boxed{\mathcal M_k(y)=
   \frac{\alpha(-1)^{k/2-1}}{\pi\sqrt2}\,
         y^{1/4}F(\sqrt y)+O_k(1).}\]\par
\subsection*{Verification of the remainder and of a counterexample sequence}
For \(T\ge1\), split the sum into \(d\le sT\) and \(d>sT\). In the first range the Bessel argument is at least \(4\pi\), and the summed remainder is \[\ll_k T^{-1/2}\sum_{s\ge1}s^{-5/2}
                      \sum_{d\le sT}d^{-1/2}
 \ll_k\sum_{s\ge1}s^{-2}\ll_k1.\] In the second range, boundedness of \(J_{k-1}\) on \([0,4\pi]\) bounds the exact contribution by \[\ll_k T\sum_{s\ge1}s^{-1}\sum_{d>sT}d^{-2}\ll_k1.\] The corresponding proposed leading terms are bounded by \[\ll T^{1/2}\sum_{s\ge1}s^{-3/2}\sum_{d>sT}d^{-3/2}\ll1.\] This proves the uniform \(O_k(1)\) remainder in [eq:correct] (\#eq:correct). It avoids the small-argument failure of the printed expression \(O(\min\{z^{-3/2},1\})\): the leading cosine term diverges at \(z=0\), whereas \(J_{k-1}(z)\) does not.\par
To demonstrate actual incompatibility of the two asymptotic statements, let \[B=\zeta(3/2)\sum_{d\ge1}Q(d)\sqrt d\in(0,\infty).\] Absolute convergence permits dominated convergence along \(T_n=n!\). For each fixed \(d,s\), eventually \(d\mid n!\), so \[F(n!)\longrightarrow-B/\sqrt2.\] All assumptions of Theorem 3 hold with fixed \(k=2\) and \(y_n=(n!)^2\to\infty\). Formula [eq:correct] (\#eq:correct) yields \[\frac{\mathcal M_2((n!)^2)}{\sqrt{n!}}\longrightarrow
            -\frac{\alpha B}{2\pi},\] whereas the printed formula would yield \(-\alpha B/\sqrt\pi\). These limits differ since \(\alpha B>0\). The contradiction is analytic and does not depend on numerical Bessel evaluation or finite sampling.\par
\subsection*{Counterevidence, impact, and verified repair}
A normalization change cannot reconcile the two displays while preserving their common \(\alpha\), standard \(J_\nu\), and common argument. The phase in the theorem statement is correct; the proof page’s \((-1)^{k/2}\) must instead be \((-1)^{k/2-1}\). The replacement [eq:correct] (\#eq:correct) and the split-range proof above repair this asymptotic step, conditional only on the checked Bessel representation. This is Major under the specified calibration because it changes a main theorem’s leading asymptotic. The demonstrated consequence is the normalization error. The order \(y^{1/4}\), universal oscillatory series, and qualitative parity dependence survive. It does not refute the original modular-form correlation result or certify the computational sign-change claim.\par
Problem 2 (An impossible remainder in Lemma 7.3: known issue reproduced).  The remainder \(R(T)\ll_\varepsilon T^{-2+\varepsilon}\) is incompatible with RH. Theorem 6’s printed proof depends on that remainder.\par
\subsection*{Source and noncancellation proof}
Equation reference: eq:dirichlet. Lemma 7.3 appears at PDF p. 46 (printed p. 672); its contour argument and Euler product are on PDF p. 47. Write \(B_p=2p^2+p-1\). The exact local identity \(A_p+B_p=p^4\) gives \[
 L(s)=\sum_{d\ge1}Q(d)d^{-s}
     =\frac{\zeta(s+2)}{\zeta(2s+4)}H(s),\qquad
 H(s)=\prod_p\left(1+\frac{B_p}{A_p(1+p^{s+2})}\right).\] Indeed, multiply one \(H\)-factor by \(1+p^{-s-2}\). The result is \(1+p^{2-s}/A_p\), the local factor of \(L(s)\).\par
On every compact subset of \(\Re s>-2\), the correction in each \(H\)-factor is \(O(p^{-\sigma_0-4})\) for some \(\sigma_0>-2\). Thus the product is locally uniformly absolutely convergent and holomorphic there. Its denominators are nonzero since \(|p^{s+2}|>1\). Assume RH and choose any nontrivial zeta zero \(\rho=\frac12+i\gamma\). At \[s_0=(\rho-4)/2=-\frac74+\frac{i\gamma}{2},\] the denominator in [eq:dirichlet] (\#eq:dirichlet) vanishes. The numerator is \(\zeta(\rho/2)\), nonzero under RH because its real part is \(1/4\). An \(H\)-factor could vanish only if \[A_p p^{s_0+2}+p^4=0.\] Taking absolute values and fourth powers would give \(A_p^4=p^{15}\), contrary to \(A_p\equiv1\pmod p\). Every local factor is nonzero at \(s_0\); absolute convergence proves \(H(s_0)\ne0\). Therefore \(L(s)\) has a genuine pole there, without any assumption that \(\rho\) is simple.\par
\subsection*{Why the asserted remainder is false}
Equation reference: eq:abel. Let \(A(x)=\sum_{d\le x}Q(d)\). The asserted estimate would give \(A(x)=L(0)-\Delta/x+R(x)\). For \(\Re s>0\), partial summation then implies \[
 L(s)=L(0)-\frac{s\Delta}{s+1}
             +s\int_1^\infty R(x)x^{-s-1}\,dx.\] Fix \(0<\varepsilon<1/4\). If \(R(x)=O(x^{-2+\varepsilon})\), the integral is holomorphic on \(\Re s>-2+\varepsilon\). Formula [eq:abel] (\#eq:abel) would continue \(L\) into that half-plane with no possible pole except \(s=-1\). But \(s_0\) belongs to that half-plane and is a nonremovable pole, contradicting uniqueness of meromorphic continuation.\par
An integer-only interpretation of \(T\) does not help: replacing \(x\) by \(\lfloor x\rfloor\) changes \(\Delta/x\) by \(O(x^{-2})\). Likewise \(d<T\) versus \(d\le T\) changes an endpoint by \(Q(T)=O(T^{-2})\). The argument holds for any fixed \(\Delta\), so the variable misprint in its displayed Euler product is immaterial.\par
\subsection*{Dependence, prior correction, and supported judgment}
Section 7.3, PDF p. 49 (printed p. 675), explicitly uses \(E(n)=n^2R(n)\ll_\delta n^\delta\) for arbitrarily small \(\delta>0\). This is the failed estimate in another form. The claim that the contour crosses only \(s=0,-1\) omits the poles at \((\rho-4)/2\). The original derivation has a Major core dependency defect.\par
This is a known-issue reproduction, not a claim of priority. Nelson, Ross, Wassercug, and Xue [nrwx], Introduction and Sections 4–5, identify this contour obstruction and state an unconditional alternative proof of the sharp-average limit, with a power-saving remainder. They also address the separate Chebyshev-weight problem in the original trace-formula argument. This review verifies the original obstruction and records the replacement source; it does not certify all steps of that external proof. The \(1/2\) limit in Theorem 6 is not disproved. No unproved replacement bound for \(R(T)\) is proposed.\par
Problem 3 (Cutoff selection does not yield the claimed bound).  The parameter after equation (7) does not balance the displayed terms; a different parameter recovers the claimed local bound.\par
\subsection*{Source, calculation, and correction}
At PDF pp. 13–14 (printed pp. 639–640), the error is bounded, up to arbitrarily small powers, by \[\frac{Y}{P}+\frac{Y}{\sqrt T}
       +X^{1/2}T^{3/16}+\frac{Y^2}{X}
       +\frac{Y(PX)^{1/2}}{T}.\] The text chooses \(T=Y^{5/6}P^{5/12}X^{-1/12}\) and says the third and last terms are equal. In fact equality requires \[T^{19/16}=YP^{1/2},\qquad \boxed{T=Y^{16/19}P^{8/19}.}\] The defect reaches the printed derivation within Theorem 1’s own parameter range. Set \[\delta_1=\frac1{100},\qquad \delta_2=\frac14,\qquad
 X=P\to\infty\ \text{through primes},\qquad Y=X^{3/4}.\] These satisfy the conditions on PDF pp. 1–2: \(0<\delta_1<1/11\), \(2\delta_1<\delta_2<(4-18\delta_1)/13\), and \(P\ll X^{1+\delta_1}\). The printed cutoff is \(T=X^{23/24}\). After the prescribed normalization \(Y^{-1}\sqrt{P/X}\), the last term of (7) becomes \(P/T=X^{1/24}\); it does not even provide a decaying error. Section 3.4 (PDF pp. 38–39) explicitly substitutes Proposition 3.1 into the proof of Theorem 1 and asserts a positive power saving throughout that parameter range. The corrected cutoff is \(T=X^{20/19}\), for which the third and last normalized terms both equal \(X^{-1/19}\). This comparison concerns the available upper bound, not a lower bound for the actual class-number error or a counterexample to the theorem.\par
With the corrected choice, the balanced normalized terms equal \(P^{11/19}Y^{-16/19}\). The second normalized term satisfies \[P^{11/38}Y^{-8/19}X^{-1/2}\le P^{11/19}Y^{-16/19}\] for \(Y<X\) and \(P>2\): their ratio is \(Y^{8/19}/(X^{1/2}P^{11/38})<1\). The remaining terms are those in Proposition 3.1. This is a Major defect in the printed core derivation, with a verified local repair using the same error estimate. The direct correction does not justify downgrading an invalid step actually used for a principal result. No failure of the conclusions of Proposition 3.1 or Theorem 1 is demonstrated, and their other proof dependencies are not certified here.\par
\section*{Other checked evidence and limits}
The Euler-factor cancellations in the Poisson/Bessel reduction were checked against PDF pp. 42–43. The companion TeX contains the same normalization, contour, and cutoff problems; this is a local comparison.\par
The recovered Section 8 calculation evaluates three sets of 15,015 sample values using a reviewer-written Hurwitz-zeta representation. Its stored extrema are negative at the integer and half-integer sample locations and positive at offset \(0.162\). That is numerical corroboration, not a rigorous certificate: the record lacks interval bounds and independent verification of the Euler-product tail. The original Mathematica notebook was unavailable in the checked materials. This absence is not itself a mathematical defect, and Corollary 4 is not labeled false on that basis. The numerical run was not repeated or misattributed to this continuation.\par
Theorem 5’s normalized smoothing presupposes \(\int_0^\infty\Phi(u)u\,du\ne0\). Its arbitrary complex-weight wording needs that domain condition; this clarification is not promoted to a substantive finding.\par
\section*{Conclusion and disposition}
The supported overall impact for the acquired version is Major, supported by the leading-order error, the invalid remainder underlying the printed sharp-average proof, and the incorrect cutoff in the derivation of Theorem 1. This is compatible with the principal phenomenon and a subsequent repair of the sharp-average result.\par
Item | Judgment | Repair verification | Author adoption | Remaining work\par
P1 | Major | Correct coefficient and bounded remainder proved. | Not documented | Full theorem proof outside scope.\par
P2 | Major; known issue | External replacement located; full proof not certified here. | Not documented | External proof verification outside scope.\par
P3 | Major | Correct cutoff verified from the displayed bound. | Not documented | Full theorem proof outside scope.\par
Cor. 4 | Needs Evidence | Numerical corroboration only. | Not applicable | Rigorous numerical certificate unverified.\par
Scientific adjudication remains pending. The independent focused checks and current build/page-inspection bindings are recorded in provenance. No human review or author adoption is asserted.\par
9 N. Zubrilina, Murmurations, Invent. Math. 241 (2025), 627–680. https://doi.org/10.1007/s00222-025-01347-8. Reviewed publisher PDF; companion https://arxiv.org/abs/2310.07681v2. NIST Digital Library of Mathematical Functions, §10.17, especially equations 10.17.2–10.17.3. https://dlmf.nist.gov/10.17. Consulted in the predecessor continuation. T. Nelson, E. Ross, M. Wassercug, and H. Xue, On murmurations, arXiv:2609.29631v1 (2026), Introduction and Sections 4–5. https://arxiv.org/abs/2609.29631v1.\par
\section*{Supporting checks}
Primary PDF SHA-256:\par
> `72dc0df9cfa2c6a09d6dbdfe5398653eda2099f9577f3bd6bb7f0b87b4dc0efa`\par
The decisive arguments in Problems 1–2 are analytic proofs. The recovered symbolic and rational checks verify only the stated algebra and exponent manipulations; numerical illustrations are not proofs.\par
The predecessor report attributes targeted Bessel and meromorphic-continuation checks to two native assistants and procedural checks to a third; the user confirms that all three tasks finished. Their separate task transcripts were not available in the preserved files, so no new reading of those transcripts is asserted.\par
In the predecessor recovery, The independent check checked the cited source evidence, retained the established arguments, completed the verification records, and performed the current compilation and page inspections. Current compilation and actual final-page inspections are recorded against the revised artifacts. Human reviews remain empty; scientific adjudication is pending.\par
\end{document}
